\documentclass[10pt,a4paper]{amsart}
\usepackage{amsmath,amssymb,amsthm,graphicx}
\usepackage{verbatim}
\newenvironment{reference}{\comment}{\endcomment}
%\newenvironment{reference}{}{}
\newenvironment{slogan}{\comment}{\endcomment}
\newenvironment{history}{\comment}{\endcomment}

% For commutative diagrams we use Xy-pic
\usepackage[all]{xy}

% We use 2cell for 2-commutative diagrams.
\xyoption{2cell}
% Theorem environments.
%
\theoremstyle{plain}
\newtheorem{theorem}[subsection]{Theorem}
\newtheorem{proposition}[subsection]{Proposition}
\newtheorem{lemma}[subsection]{Lemma}

\theoremstyle{definition}
\newtheorem{definition}[subsection]{Definition}
\newtheorem{example}[subsection]{Example}
\newtheorem{exercise}[subsection]{Exercise}
\newtheorem{situation}[subsection]{Situation}

\theoremstyle{remark}
\newtheorem{remark}[subsection]{Remark}
\newtheorem{remarks}[subsection]{Remarks}

\numberwithin{equation}{subsection}

% Macros
%
\def\lim{\mathop{\mathrm{lim}}\nolimits}
\def\colim{\mathop{\mathrm{colim}}\nolimits}
\def\Spec{\mathop{\mathrm{Spec}}}
\def\Hom{\mathop{\mathrm{Hom}}\nolimits}
\def\Ext{\mathop{\mathrm{Ext}}\nolimits}
\def\SheafHom{\mathop{\mathcal{H}\!\mathit{om}}\nolimits}
\def\SheafExt{\mathop{\mathcal{E}\!\mathit{xt}}\nolimits}
\def\Sch{\mathit{Sch}}
\def\Mor{\mathop{\mathrm{Mor}}\nolimits}
\def\Ob{\mathop{\mathrm{Ob}}\nolimits}
\def\Sh{\mathop{\mathit{Sh}}\nolimits}
\def\NL{\mathop{N\!L}\nolimits}
\def\CH{\mathop{\mathrm{CH}}\nolimits}
\def\proetale{{pro\text{-}\acute{e}tale}}
\def\etale{{\acute{e}tale}}
\def\QCoh{\mathit{QCoh}}
\def\Ker{\mathop{\mathrm{Ker}}}
\def\Im{\mathop{\mathrm{Im}}}
\def\Coker{\mathop{\mathrm{Coker}}}
\def\Coim{\mathop{\mathrm{Coim}}}

% Boxtimes
%
\DeclareMathSymbol{\boxtimes}{\mathbin}{AMSa}{"02}

%
% Macros for moduli stacks/spaces
%
\def\QCohstack{\mathcal{QC}\!\mathit{oh}}
\def\Cohstack{\mathcal{C}\!\mathit{oh}}
\def\Spacesstack{\mathcal{S}\!\mathit{paces}}
\def\Quotfunctor{\mathrm{Quot}}
\def\Hilbfunctor{\mathrm{Hilb}}
\def\Curvesstack{\mathcal{C}\!\mathit{urves}}
\def\Polarizedstack{\mathcal{P}\!\mathit{olarized}}
\def\Complexesstack{\mathcal{C}\!\mathit{omplexes}}
% \Pic is the operator that assigns to X its picard group, usage \Pic(X)
% \Picardstack_{X/B} denotes the Picard stack of X over B
% \Picardfunctor_{X/B} denotes the Picard functor of X over B
\def\Pic{\mathop{\mathrm{Pic}}\nolimits}
\def\Picardstack{\mathcal{P}\!\mathit{ic}}
\def\Picardfunctor{\mathrm{Pic}}
\def\Deformationcategory{\mathcal{D}\!\mathit{ef}}
\usepackage{cleveref}
\crefname{originalreference}{original reference}{original references}
\Crefname{originalreference}{Original reference}{Original references}
\makeatletter
\newlabel{situation-finitely-generated}{{09H6}{source}{}{sourcepointer-1}{}}
\newlabel{situation-finitely-generated@cref}{{[originalreference][0][]09H6}{[1][0][]source}{}{}{}}
\makeatother
\title{F-embeddings of a finitely generated extension correspond to compatible root tuples}
\author{The Stacks Project Authors}
\date{stacks-project:fields.tex:local-snapshot-20261012; GNU FDL 1.2 or later}
\begin{document}
\maketitle
\noindent\textit{Editorial note.} All delivered mathematical text is the author's own, copied from the pinned Stacks Project chapter file fields.tex; no statement or proof step was rewritten, reordered or removed.\par\medskip
\noindent\textit{Editorial note.} The counted proof is the author's complete argument: This bijection is the input for the product formula |Mor\_F(K, overline\{F\})| = prod deg\_s(P\_i) (b8 sibling card L48380).\par\medskip
\noindent\textit{Editorial note.} The label lemma-count-embeddings names the counted statement, and the proof environment immediately following it in the source is delivered verbatim as its proof.\par\medskip
\noindent\textit{Editorial note.} The counted statement is carried by the author's own lemma environment in the chapter source.\par\medskip
\noindent\textit{Allowed standard prerequisites (editorial).}\begin{itemize}
\item Chapter Situation situation-finitely-generated (source lines 1485-1496): K = F(alpha\_1..alpha\_n) finite, tower K\_i = F(alpha\_1..alpha\_i), P\_i the minimal polynomial of alpha\_i over K\_\{i-1\}, algebraic closure overline\{F\} fixed
\item Universal property of quotient rings: a homomorphism out of K\_\{i-1\}[x]/(P\_i) over K\_\{i-1\} is the same as a choice of a root of P\_i\textasciicircum{}\{phi\_\{i-1\}\} in the target; this is also why each step map phi\_i exists and is unique
\item The section paragraph immediately preceding the lemma (source lines 1499-1505): a morphism phi: K -> overline\{F\} yields step maps phi\_i and phi(alpha\_i) is a root of the conjugated polynomial P\_i\textasciicircum{}\{phi\_\{i-1\}\}
\end{itemize}\medskip
\ifdefined\hypertarget\hypertarget{sourcepointer-1}{}\fi
\noindent\textit{Source reference.} Original reference 09H6: sources/stacks-fields/pinned-original-source.tex, source lines 1486--1486. Prerequisite pointer invoked at 1509; it names a standard fact of the same chapter that the counted argument uses, and every step of the counted proof built on it is retained verbatim in the delivered body.\par\medskip
\begin{lemma}
\label{lemma-count-embeddings}
In Situation \ref{situation-finitely-generated} the correspondence
$$
\Mor_F(K, \overline{F})
\longrightarrow
\{(\beta_1, \ldots, \beta_n)\text{ as below}\},
\quad
\varphi \longmapsto (\varphi(\alpha_1), \ldots, \varphi(\alpha_n))
$$
is a bijection. Here the right hand side is the set of $n$-tuples
$(\beta_1, \ldots, \beta_n)$ of elements of $\overline{F}$ having
the following property:
\begin{enumerate}
\item $\beta_1 \in \overline{F}$ is a root of $P_1$;
let $\varphi_1 : K_1 \to \overline{F}$ be the homomorphism
over $F$ sending $\alpha_1$ to $\beta_1$,
\item $\beta_2 \in \overline{F}$ is a root of $P_2^{\varphi_1}$;
let $\varphi_2 : K_2 \to \overline{F}$ be the homomorphism
extending $\varphi_1$ sending $\alpha_2$ to $\beta_2$,
\item and so on until,
\item $\beta_n \in \overline{F}$ is a root of $P_n^{\varphi_{n - 1}}$.
\end{enumerate}
In each step the homorphism $\varphi_i$ exists and is unique because
$K_i = K_{i - 1}[x]/(P_i)$ and $\beta_i$ is a root of $P_i^{\varphi_{i - 1}}$.
\end{lemma}

\begin{proof}
The map from left to right is discussed above the lemma.
Let $(\beta_1, \ldots, \beta_n)$ be an element of the right hand side.
Then we let $\varphi : K = K_n \to \overline{F}$ be the unique homorphism
extending $\varphi_{n - 1}$ sending $\alpha_n$ to $\beta_n$.
Uniqueness implies that the two constructions are mutually inverse.
\end{proof}

\end{document}
